\section{Related Work}
\label{sec:related_work}

\subsection{Few-Shot Learning in Plant Pathology}
Foliar pathology often involves novel diseases and region-specific conditions with few labeled samples~\cite{shoaib2023advanced,sun2024few}. Metric-based few-shot methods address this constraint by mapping query images to support prototypes in a learned metric space~\cite{snell2017prototypical}. Several agricultural studies adapt this idea to crop foliage: Zhao et al.~\cite{zhao2024mafde} introduced multiscale attention (MAFDE-DN4) for localized lesions, Yang et al.~\cite{yang2024laboratory} benchmarked metric, optimization-based, and non-meta-learning variants for laboratory-to-field transfer, and Mu et al.~\cite{mu2024few} used supervised contrastive learning (SC-FSL) to separate visually subtle pathologies. Kudo et al.~\cite{kudo2024few} studied metric domain adaptation, and Babu et al.~\cite{babu2025agrifewnet} introduced AgriFewNet for foliar few-shot classification.

Standard metric approaches still rely on unweighted distances such as Euclidean distance or cosine similarity. They weight all feature dimensions uniformly and ignore directional feature dependencies~\cite{bateni2020improved}, so background clutter and variable lighting degrade predictions in field imagery~\cite{yang2024laboratory}. Optimization-based meta-learning methods such as MAML~\cite{finn2017maml} also require iterative gradient updates during adaptation, which conflicts with the memory and power budgets of agricultural edge processors~\cite{radocaj2025evolution}.

\subsection{Frozen Foundation Models and Parameter-Efficient Adaptation}
Large vision and vision--language models pretrained on broad data have shifted the few-shot paradigm from training from scratch to adapting fixed representations. Self-supervised transformers such as DINOv2~\cite{oquab2023dinov2} provide local patch features without label supervision, while multimodal models such as CLIP~\cite{radford2021clip} and BioCLIP~\cite{stevens2024bioclip} align visual features with botanical taxonomy.

Several parameter-efficient methods transfer these representations without full backbone retraining, including Low-Rank Adaptation (LoRA~\cite{hu2022lora}) and non-parametric key--value caches (Tip-Adapter~\cite{zhang2022tip}). PEFT reduces the number of updated parameters, but gradient-based methods still require backward passes, optimizer state, and autograd graphs~\cite{hu2022lora}, which conflict with the memory and power budgets of agricultural edge processors~\cite{radocaj2025evolution}. Tip-Adapter avoids gradient steps but operates on a single pooled embedding with an isotropic affinity and does not model class-conditional covariance~\cite{zhang2022tip}. Classical methods such as Fisher discriminant analysis~\cite{fisher1936use} and Ledoit--Wolf shrinkage~\cite{ledoit2004well} give numerically stable closed-form projections in the small-sample regime, yet their use for adapting frozen foundation models on edge hardware has received little attention in plant pathology.

\subsection{Transductive Few-Shot Inference}
Transductive few-shot methods use the query set as well as the support set and are strong baselines when the backbone is frozen~\cite{boudiaf2020transductive,ziko2020laplacian}. LaplacianShot~\cite{ziko2020laplacian} adds a graph Laplacian term that encourages queries with similar embeddings to share a label. TIM~\cite{boudiaf2020transductive} optimizes query logits for confident predictions with a class-balanced marginal prior. BD-CSPN~\cite{liu2020bdcspn} decomposes prototype estimation error into internal and cross-category distortion, correcting both through confident pseudo-labels and query feature alignment.

All three formulations require access to the complete set of query images simultaneously to construct pairwise affinity graphs or estimate marginal class distributions. While effective for batched offline benchmarks, this transductive batch requirement is fundamentally incompatible with sequential online field scouting, where a mobile terminal or autonomous scout inspects and diagnoses foliage one leaf at a time.

\subsection{On-Device Edge Cascades and Uncertainty Routing}
Deploying vision models directly on agricultural edge platforms requires bounded latency, low memory footprints, and low power draw~\cite{radocaj2025evolution,sharif2026real}. To balance accuracy and efficiency, cascaded architectures and early-exit networks allocate computation adaptively per query~\cite{viola2001rapid,teerapittayanon2016branchynet}, while selective classification withholds predictions when confidence is insufficient~\cite{geifman2017selective,traub2024overcoming}. BioDINO-Edge synthesizes these paradigms: confident queries exit immediately after Tier-1, whereas ambiguous leaves are escalated to multimodal Tier-2 verification rather than rejected, dynamically matching computation to diagnostic difficulty.

\subsection{Positioning BioDINO-Edge}
In plant pathology, most systems send every leaf through one static pipeline~\cite{zhao2026lightweight}, and their softmax scores are not calibrated for selective triage~\cite{guo2017calibration}. This has two costs: clear leaves spend latency and energy on a full evaluation, while ambiguous lesions can be misclassified with high confidence, leading to chemical over-application or undetected spread of pathogens. BioDINO-Edge addresses both by combining frozen foundation backbones, closed-form shrinkage projections, and uncertainty-gated routing. Instead of running the deep multimodal pipeline on every frame, it checks Tier-1 certainty with normalized Shannon entropy and top-2 margin and escalates only difficult leaves, which keeps diagnostic reliability high within tight edge budgets.
